\chapter{Conclusiones y Trabajos futuros}
En este capítulo se exponen los logros y objetivos alcanzados a lo largo del desarrollo del proyecto \Gls{snapperml}, así como las principales lecciones aprendidas durante el proceso. También se detalla el trabajo futuro que podría llevarse a cabo para seguir mejorando el \gls{framework} y consolidarlo como una herramienta robusta y eficiente dentro del ámbito \Gls{mlopsg}. Finalmente, se presenta una reflexión sobre la importancia de las prácticas adoptadas y el impacto del proyecto en la comunidad de aprendizaje automático.

\section{Logros y Objetivos Alcanzados}
El desarrollo de \Gls{snapperml} ha cumplido con los objetivos propuestos, logrando avances significativos que refuerzan la importancia de las buenas prácticas en la gestión de la \gls{deuda} y la implementación de una interfaz web intuitiva. A continuación, se resumen los principales logros:

\begin{itemize} 
    \item \textbf{Gestión y Resolución de \Gls{deuda}}: Se ha abordado la \gls{deuda} acumulada mediante una profunda refactorización del código y la optimización de componentes clave. Esto ha resultado en un sistema más eficiente, mantenible y escalable, con dependencias actualizadas y un enfoque proactivo para evitar \gls{deuda} futura. Además, se han implementado pruebas automatizadas para garantizar la estabilidad y calidad del \gls{framework}. 
    \item \textbf{Diseño e Implementación de la Interfaz Web}: Se ha desarrollado una interfaz web moderna usando React y Bootstrap, que simplifica el uso del \gls{framework} para los usuarios. La interfaz permite cargar archivos \acrshort{yaml}, configurar experimentos y ejecutar comandos de manera eficiente desde un navegador. Este avance ha hecho que \Gls{snapperml} sea más accesible y usable, especialmente para aquellos sin experiencia avanzada en la línea de comandos. 
    \item \textbf{Pruebas y Validación}: Se han implementado y ejecutado pruebas unitarias, de integración y de sistema, validando que \Gls{snapperml} cumple con los requisitos funcionales y no funcionales. El sistema ha demostrado ser capaz de manejar múltiples solicitudes concurrentes, garantizando un alto rendimiento y una experiencia de usuario fluida. Las pruebas también han asegurado la reproducibilidad y trazabilidad de los experimentos. 
\end{itemize}
Estos logros refuerzan el compromiso de \Gls{snapperml} con las prácticas modernas de desarrollo de software y establecen una base sólida para futuras mejoras.

\section{Lecciones Aprendidas} El desarrollo de \Gls{snapperml} ha sido una experiencia enriquecedora, que ha aportado numerosas lecciones importantes:
\begin{itemize} 
    \item \textbf{Importancia de la Planificación y Gestión del Tiempo}: El uso de herramientas como Kanban y Wakatime ha sido fundamental para organizar tareas y optimizar el tiempo disponible, especialmente al compaginar el trabajo académico con las prácticas profesionales. Esta experiencia ha demostrado que una buena planificación es esencial para mantener un progreso constante y cumplir con los plazos. 
    \item \textbf{Gestión de \Gls{deuda}}: La \gls{deuda} puede paralizar el desarrollo si no se aborda adecuadamente. Refactorizar el código y establecer prácticas de prevención ha mejorado la calidad general del proyecto y ha facilitado la implementación de nuevas funcionalidades. Esta lección destaca la necesidad de mantener un equilibrio entre el desarrollo rápido y la sostenibilidad a largo plazo. 
    \item \textbf{Flexibilidad y Adaptabilidad}: La integración de nuevas tecnologías como xterm.js y Vite ha requerido una mentalidad abierta y una rápida adaptación. La capacidad de aprender y aplicar soluciones innovadoras es clave en proyectos tecnológicos, especialmente cuando los requisitos evolucionan. 
\end{itemize}
Estas lecciones no solo han contribuido al éxito del proyecto, sino que también proporcionan una base de conocimiento útil para futuros desarrollos.

\section{Trabajo futuro}
El proyecto \Gls{snapperml} se encuentra en una posición prometedora, pero todavía hay espacio para mejoras y expansiones. Se sugieren varias líneas de trabajo futuro:

\begin{itemize} 
    \item \textbf{Ampliación de Funcionalidades}: Desarrollar nuevas características en la interfaz web, como la visualización de gráficos de resultados de experimentos y el soporte para la optimización de hiperparámetros. También se podría mejorar la gestión de logs y la visualización de métricas en tiempo real. 
    \item \textbf{Integración de Herramientas Avanzadas de \Gls{mlopsg}}: Investigar la integración con plataformas avanzadas como Kubeflow o Vertex AI para mejorar las capacidades de despliegue y monitorización. Estas integraciones permitirían a \Gls{snapperml} ofrecer un ciclo de vida de modelos más completo, desde el desarrollo hasta la producción. 
    \item \textbf{Automatización y Mejora de Pruebas}: Ampliar el sistema de pruebas con pruebas de carga y estrés para asegurar que \Gls{snapperml} pueda gestionar grandes volúmenes de datos y usuarios de manera eficiente. Esto también incluiría la automatización de pruebas de seguridad y la validación de configuraciones. 
    \item \textbf{Seguridad y Escalabilidad}: Implementar medidas de seguridad avanzadas, como la gestión de usuarios y permisos, para proteger los datos y garantizar el cumplimiento de las normativas de privacidad. Asimismo, optimizar la infraestructura para permitir un escalado eficiente en entornos cloud, garantizando una experiencia sin interrupciones. 
    \item \textbf{Comercialización y Plataforma en la Nube}: Explorar la posibilidad de desarrollar una versión comercial de \Gls{snapperml} como una plataforma en la nube. Esta plataforma podría ofrecer funcionalidades avanzadas como la gestión de usuarios, gobernanza de modelos y servicios de colaboración en tiempo real, con un enfoque en la reproducibilidad y las mejores prácticas de \Gls{mlopsg}. 
\end{itemize}
Estas líneas de trabajo futuro no solo mejorarían la funcionalidad y usabilidad del \gls{framework}, sino que también abrirían nuevas oportunidades para su aplicación en entornos académicos y empresariales.

\section{Conclusión Final} 
El desarrollo de \Gls{snapperml} ha sido un desafío técnico y académico que ha aportado resultados valiosos y aprendizajes significativos. La gestión efectiva de la \gls{deuda} ha transformado el \gls{framework} en una herramienta más eficiente y mantenible, mientras que la nueva interfaz web ha mejorado la accesibilidad y usabilidad para los usuarios. Este proyecto destaca la importancia de las prácticas de \Gls{mlopsg} en la gestión de experimentos de aprendizaje automático, especialmente en términos de reproducibilidad y sostenibilidad del software.

Con las mejoras realizadas y las propuestas de trabajo futuro, \Gls{snapperml} está preparado para seguir evolucionando y consolidarse como una herramienta esencial para la comunidad de aprendizaje automático. Este marco de trabajo no solo ayuda a investigadores y científicos de datos a gestionar sus experimentos de manera más eficiente, sino que también promueve el uso de buenas prácticas en la industria. En definitiva, \Gls{snapperml} representa un avance significativo hacia la creación de sistemas de aprendizaje automático más robustos, escalables y accesibles.